\chapter{Finite algebra and normalization identities}\label{app:finite}

This appendix proves the finite identities on which the asymptotic argument rests. The first calculation compares the chosen crossing matrices and reduced trace with the conventional invariant. The next two calculations establish the Newton basis, its genuine geometric turn, and the exact color decomposition. The last calculation expands an arbitrary twist power while keeping derivatives at repeated roots. No limiting argument is used in these identities.

\section{The crossing comparison and enhancement}\label{app:normalization}
This section supplies the finite normalization calculation used in Proposition~\ref{prop:normalization}. Put $\omega=Q$ and
\[
 p_j=\prod_{h=1}^{j}(1-\omega^h),\qquad
 \bar p_j=\prod_{h=1}^{j}(1-\omega^{-h}),\qquad p_0=1.
\]
Kashaev's unbraided four-index matrix may be written by introducing nonnegative residues
\[
 u=[j-i-1],\quad v=[i-l],\quad t=[l-k],\quad h=[k-j].
\]
Its support conditions are equivalent to $u+v+t+h=N-1$, and on that support
\begin{equation}\label{eq:A-original-crossing}
 R^o_{(i,j);(k,l)}=
 \frac{N\omega^{1-(t+h)(u+h)}}{p_u\bar p_vp_t\bar p_h}.
\end{equation}
One free $j$ and such a composition specify one entry. Let $F_{ij}=\omega^{ij}/\sqrt N$. Its full Fourier transform has only the ordinary-integer-conserving entries
\begin{equation}\label{eq:A-Fourier-crossing}
 ((F^{-1}\otimes F^{-1})R^o(F\otimes F))_{(a-k,b+k);(a,b)}
 =\omega^{1+(a-k)(b+1)}\frac{p_ap_{b+k}}{p_{a-k}p_bp_k},
\end{equation}
for $0\le k\le\min(a,N-1-b)$.

Here is the finite summation. Define
\[
 S_m(X)=\sum_{j=0}^m\frac{X^j}{p_j\bar p_{m-j}}
 =\frac{\prod_{j=1}^m(X-\omega^j)}{p_m}.
\]
The equality follows by induction from $\bar p_m=(-1)^m\omega^{-m(m+1)/2}p_m$ and
$(1-\omega^m)S_m=(X-\omega^m)S_{m-1}$. The free Fourier index imposes conservation modulo $N$. The two remaining factors are
$S_{N-1-z}(\omega^{c-z})S_z(\omega^{z-a})$; their zeros force $c\le z\le a$. With $m=a-c$, the residual sum is
\[
 \omega^{1+c(b+1)}\frac{p_a}{p_c}
 \sum_{j=0}^m\frac{(-1)^j\omega^{j(j-1)/2}\omega^{(b+1)j}}
 {p_jp_{m-j}}
 =\omega^{1+c(b+1)}\frac{p_a}{p_c}
 \frac{(\omega^{b+1};\omega)_m}{p_m}.
\]
If $b+m\ge N$, the product contains $1-\omega^N$ and vanishes. Otherwise it is $p_{b+m}/p_b$, proving \eqref{eq:A-Fourier-crossing} with integer, rather than merely residue, conservation.

Set $G_j=p_j/(1-\omega)^j=q^{j(j-1)/2}[j]_q!$, $D_j=(-q)^j$, and $\gamma=\omega q^{-d^2/2}$. Substitution into \eqref{eq:universal-crossing} gives exactly
\[
 (F^{-1}\otimes F^{-1})R^o(F\otimes F)
 =\gamma(DG\otimes G)\mathscr R(G^{-1}\otimes G^{-1}D^{-1}).
\]
For the enhanced braid conventions take
\[
 B_K=P(R^o)^T,\quad\mu_K=-qS,\quad\alpha_K=-q,
 \qquad B_q=P\mathscr R^T,\quad\mu_q=K^{-1},\quad
 \alpha_q=q^{(N^2-1)/2},
\]
where $S v_j=v_{j+1\bmod N}$. On $b$ strands, the single similarity
\[
 U_b=F^{-\otimes b}\bigotimes_{j=1}^b(G^{-1}D^{j-1})
\]
works for every adjacent generator because the quantum crossing conserves integer weight and commutes with $D\otimes D$. It also sends $\mu_q^{\otimes b}$ to $\mu_K^{\otimes b}$, and $\gamma\alpha_q=\alpha_K$. The writhe characters cancel exactly. The partial traces commute with the remaining module generators; distinct $K$ weights and nonzero $F$ arrows imply scalarity. Thus every internal similarity cancels, including the one open factor of a reduced cut. This proves the finite enhanced comparison used in Chapter~\ref{ch:source}. Large factorial gauges are never estimated by their norms.

The primal/dual change of basis is $Bv_j=(-1)^jq^{j(j-N)}v_{d-j}$. Direct generator substitution conjugates $K^{-1}$ to the primal pivot and the crossing to the recorded primal convention. At completeness $g_{\rm ADO}=(-1)^dK$ and the twist has the corresponding sign. A $b$-strand closure of writhe $w$ and $c$ components has $w\equiv b-c\pmod2$, so the remaining enhancement ratio is $(-1)^{d(c-1)}$. Its modulus is one, and for a knot it is one in this convention.

\section{The block polynomial and the physical Newton turn}\label{app:newton}
For completeness we justify the block minimal polynomial used in Lemma~\ref{lem:commutant}. Before specialization, the ordinary highest-weight module of highest weight $d$ has weights $d-2a$. In the pair, a highest vector of total degree $j$ is obtained by solving the bidiagonal raising equation. If its coefficient at $|a,j-a\rangle$ is $u_a$, the equation at $|a,j-a-1\rangle$ determines $u_{a+1}/u_a$ as the negative ratio of the two nonzero raising entries. Thus for generic $q$ there is one such vector in each degree $j=0,\ldots,d$. Its highest weight is $2d-2j$, its Casimir is $q^{2d-2j+1}+q^{-2d+2j-1}$, and its successive lowering vectors are nonzero through the usual length. The dimensions $2d-2j+1$ sum to $N^2$. Distinct generic Casimirs make these submodules independent.

In degree $T\le d$, precisely the highest vectors $j\le T$ contribute one descendant each. Their eigenvalues give the block characteristic polynomial. At $q=e^{\pi\ii/N}$ it becomes
\[
 \prod_{j=0}^{T}(z-c_j),\qquad c_j=q^{2j+1}+q^{-2j-1}.
\]
All physical matrix entries specialize regularly. The nonzero adjacent entries of the tridiagonal block make a boundary vector cyclic, so this characteristic polynomial is also minimal even when roots coincide. At $T=d$ it is $P_N$. The elementary identity
\[
 \prod_{j=0}^{N-1}(z-q^{2j+1}-q^{-2j-1})=2T_N(z/2)+2
\]
can also be checked by putting $z=w+w^{-1}$: the roots correspond to $w^N=-1$, and the monic degree and leading coefficient agree. The generating property of $W_d$ propagates its vanishing to the full pair.

Now $B_r(\Omega)$ annihilates every degree block $W_T$ with $T<r$, with multiplicities, and has first-color bandwidth at most $r$. In the turned matrix its first-color displacement is $d-(a+b)$, so $\Turn B_r$ annihilates every $W_T$ with $T<d-r$. Its representing polynomial is therefore divisible by $B_{d-r}$.

For the reverse degree bound, its first middle-block column is
\[
 (\Turn B_r)_{a0}=x^{a-d}(B_r(\Omega))_{(d-a,0);(d-a,0)},
\]
which vanishes for $a>d-r$. A polynomial of degree $k<N$ in the cyclic tridiagonal middle block has a nonzero $k$th coordinate on the first basis vector, from its unique outermost path. Thus $\deg\Turn B_r\le d-r$, proving $\Turn B_r=C_rB_{d-r}$.

Use the physical entry from $(d-r,0)$ to $(0,d-r)$. On the turned side it is the unique $r$-step path in the high-degree block, whose factors are
\[
 Q^{j-r}(1-Q^{r-j})(1-Q^{j+1}),\qquad0\le j<r.
\]
Including the leading coefficient $(-1)^r$ gives $s_r^2$. The corresponding path for $B_{d-r}$ in the low block gives $s_{d-r}^2$. All turn monomials in this entry equal one. Hence
\[
 \Turn B_r=\frac{s_r^2}{s_{d-r}^2}B_{d-r},
 \qquad \Turn\Bhat_r=\Bhat_{d-r}.
\]
Finally $s_rs_{d-r}=N$ and \eqref{eq:leading-trace} give
$\phi(\Bhat_i\Bhat_j)=\delta_{i+j,d}$: use degree below $d$, leading coefficient at $d$, and full divisibility by $P_N$ above it. This proves Proposition~\ref{prop:newton-turn} by physical support and exact paths.

\section{Circle coefficients and the signed alias}\label{app:colors}

We derive the coefficient formulas rather than infer them from values at the distinct Casimir roots. That distinction is necessary because the quotient algebra also contains first derivatives at double roots. The generating-function calculation first gives a polynomial identity at a generic parameter; regular specialization then retains all those derivatives automatically.
The generic Chebyshev--Newton identity is
\begin{equation}\label{eq:A-Cheb-newton}
 U_{j-1}(z/2)=\sum_{r=0}^{j-1}(-1)^r
 \genfrac{[}{]}{0pt}{}{j+r}{2r+1}_q B_r(z).
\end{equation}
To prove it, successively divide
$1-zt+t^2=(1-c_rt+t^2)+t(c_r-z)$. The formal identity is
\[
 \frac1{1-zt+t^2}
 =\sum_{r\ge0}\frac{(-t)^rB_r(z)}{\prod_{i=0}^{r}(1-c_it+t^2)}.
\]
After $r$ steps the remainder is divisible by $t^{r+1}$. The denominator factors as a consecutive $q$-string, and the Gaussian Pascal recurrence gives its reciprocal coefficient. Taking the coefficient of $t^{j-1}$ gives \eqref{eq:A-Cheb-newton}. For $j\le\lfloor N/2\rfloor$, every factorial index is below $N$; specializing and turning proves \eqref{eq:low-alias}.

For higher colors use the polynomial identity
\[
 U_{N-h-1}=\tfrac12(\Cheb_hU_{N-1}-\Cheb_NU_{h-1}).
\]
Since $\Cheb_N=-2$ in the full algebra,
\[
 L_{N-h}=(-1)^h\tfrac12\Cheb_hL_N+(-1)^NL_h.
\]
The $v$th turned coefficient of the first term is
$(-1)^hN\sum_{k=0}^{d}\cos(h(2k+1)\pi/N)\Bhat_v(c_k)$.
This finite sum can be evaluated without a zero denominator. Put $z=w+w^{-1}$. The factorization of $B_v$ is the consecutive string
\[
 B_v(w+w^{-1})=(-1)^vw^{-v}
 \prod_{j=0}^{2v-1}\left(w-q^{-(2v-1)+2j}\right).
\]
For $0\le v,h\le M$, all its Fourier exponents after multiplication by $w^{\pm h}$ have absolute value below $N$. Summing at $w=q^{2k+1}$ extracts only the constant coefficient. The finite Gaussian product gives that coefficient as zero for $v<h$, and as
$(-1)^h s_{2v}/(s_{v-h}s_{v+h})$ for $v\ge h$. Dividing by $s_v^2$ and multiplying the two factors of $N$ proves \eqref{eq:core-alias}, including $h=0$ and both parities.

For the half twist, multiplication by $z$ in Newton coordinates is the lower bidiagonal matrix $Z$ with $Z_{rr}=c_r$ and $Z_{r+1,r}=-4\sin^2((r+1)\pi/N)$. The one-crossing matrix is
\[
 (R(Z))_{ij}=\kappa(-1)^iq^{i(j+1)}\frac{s_i}{s_js_{i-j}},\qquad i\ge j.
\]
Its first column is the finite Newton expansion of the crossing. Substituting the two neighboring ratios shows that it commutes with $Z$; since the first coordinate is cyclic, this determines the full matrix. Reversal on both sides gives \eqref{eq:half-twist}. Coefficient conjugation gives the inverse crossing because its full Hermite values and logarithmic derivatives are conjugate. This argument retains all normal directions.

\section{Fourier--Hermite interpolation of an arbitrary twist}\label{app:fourier}

Ordinary interpolation prescribes values. Hermite interpolation also prescribes derivatives at a repeated node. Here the crossing power supplies both data, and the change of variable $z=w+w^{-1}$ turns the interpolation problem into a finite Fourier calculation. The two Fourier equations below solve simultaneously for the coefficients that ordinary value interpolation would fail to distinguish.
For an integer $\ell$, let $f_{N,\ell}$ be the degree-below-$N$ remainder of $r_N^\ell$ modulo $P_N$. Negative powers exist because all its nodal values are nonzero. At $w_k=q^{2k+1}$ define
\[
 h_k=\bigl((-1)^kq^{k(k+1)}\bigr)^\ell,
 \qquad d_k=-\frac\ell2(d-2k)h_k.
\]
At a double node, $d_k=(w_k-w_k^{-1})f_{N,\ell}'(c_k)$; at the odd simple node both sides of the logarithmic derivative condition vanish without prescribing an extra derivative. Put
\[
 H_n=\frac1N\sum_kh_kw_k^{-n},\qquad
 D_n=\frac1N\sum_kd_kw_k^{-n},\qquad0\le n<N.
\]
Then
\begin{equation}\label{eq:A-Hermite-Fourier}
 f_{N,\ell}(z)=a_0+\sum_{n=1}^{d}a_n\Cheb_n(z),\qquad
 a_0=H_0,\quad a_n=\frac{D_n+(N-n)H_n}{N}.
\end{equation}
Indeed the Laurent polynomial $f(w+w^{-1})$ has symmetric coefficients $a_n$ at $w^{\pm n}$. At $w_k^N=-1$ its values and logarithmic derivatives have aliases
\[
 H_n=a_n-a_{N-n},\qquad
 D_n=na_n+(N-n)a_{N-n}.
\]
Solving gives \eqref{eq:A-Hermite-Fourier}. At the even middle index $H_{N/2}=0$ and the derivative equation still determines the coefficient. Paired nodes give the required compatibility, and the odd simple root imposes exactly one condition. Thus no root or normal is omitted.

Set $a_N=a_{N+1}=0$ and
\[
 \beta_j^{\ell/2}=(-1)^j(a_j-a_{j+2}),\qquad0\le j<N.
\]
The identity $\Cheb_n=U_n-U_{n-2}$ gives the circle expansion. Since $|h_k|=1$ and $|d_k|\le|\ell|N/2$, one has $|a_0|\le1$, $|a_n|\le1+|\ell|/2$, and
\[
 \sum_j|\beta_j^{\ell/2}|\le2N(1+|\ell|/2).
\]
For $\ell=2P$, these are the coefficients used in \eqref{eq:sampled-derivative}. To obtain that formula directly, substitute the two finite Fourier sums into $a_j-a_{j+2}$, set $j=N-h-1$, and pair $h$ with $-h$ in an even corridor. The difference of the two exponentials gives $\sinh(2\pi\ii z_k)$; the differentiated data give the linear factor $2PNz_k+h$, and differentiating the hyperbolic sine gives its cosine. Their sum is precisely $U_N'(z_k)/(2\pi\ii)$. The constant phase is $(-1)^d\kappa^{-P}$ and the two averaging factors give $N^{-2}$. This verifies \eqref{eq:sampled-derivative} with both normal terms and both parity conventions.
